\documentclass[../../main.tex]{subfiles}
\begin{document}

\begin{enumerate}[(1)]
  \item$\vec{n_0}$は互いに反対を向いた$2$つのベクトルである．$S_0$の法線ベクトルの一つに
    $\vec{n}=(1,2,-2)$があるから，求めるベクトルは$k\in\mathbb{R}$として$k\vec{n}$の形で書ける．この絶対値が$1$であるから，$k=\pm\dfrac{1}{3}$.故に
    $\vec{n_0}=\pm\dfrac{1}{3}(1,2,-2)\cdots\text{(答)}$が求める答えである．

  \item$l$の方向ベクトル$\vec{l}=(-1,1,0)$である．又，$\vec{n_0}=$を$\vec{l}$の周りに回転させたベクトルを考えればよい．$l$上に点$A$をとる．以下位置ベクトルの基準点を$A$として，点
\end{enumerate}
\newpage
\end{document}
